\subsection{等变 Donaldson–Fukaya 范畴}\label{sec:14b}
值得看看上述定义如何融入小节 \sr{sec:10b}{(10b)} 与 \sr{sec:12i}{(12i)} 中较抽象的讨论。

假设给定带协调 $\Z/2$ 作用的范畴 $A$。该作用由自等价 $E:A\to A$ 和函子同构 $T:E^2\to\Id_A$ 组成，满足唯一的协调条件 $R_E(T)=L_E(T)$，更具体地，
\begin{equation}\label{eq:14.6}
T_{E(X)}=E(T_X)\in\Hom_A(E^3(X),E(X)),\qquad\text{对所有 }X.
\end{equation}
于是可以如下定义新范畴 $A^{\mathrm{strict}}$。对象形如 $(X_\pm,\allowbreak F_+,\allowbreak F_-)$，其中 $X_\pm$ 为 $A$ 的对象，$F_\pm:X_\pm\to E(X_\mp)$ 为同构，使 $E(F_-)\circ F_+:X_+\to E^2(X_+)$ 等于 $T_{X_+}^{-1}$。由\eref{14.6}，从 $X_-$ 出发的相应性质自动成立。$A^{\mathrm{strict}}$ 中从 $(X_{0,\allowbreak \pm},\allowbreak F_{0,\allowbreak \pm})$ 到 $(X_{1,\allowbreak \pm},\allowbreak F_{1,\allowbreak \pm})$ 的态射，就是 $A$ 中从 $X_{0,\allowbreak +}$ 到 $X_{1,\allowbreak +}$ 的态射；立即可见 $A^{\mathrm{strict}}$ 等价于 $A$。不过，现在有了严格 $\Z/2$ 作用：函子 $E^{\mathrm{strict}}$ 交换对象的正、负两部分，并将
\[
\begin{aligned}
\Phi\in\Hom_A(X_{0,+},X_{1,+})\longmapsto{}
F_{1,-}^{-1}\circ E(\Phi)\circ F_{0,-}\in\Hom_A(X_{0,-},X_{1,-}).
\end{aligned}
\]
特别地，$A^{\mathrm{strict}}$ 中在 $E^{\mathrm{strict}}$ 下不变的对象，由 $A$ 的单个对象 $X$ 及同构 $F:X\to E(X)$ 给出，其中 $E(F)\circ F=T_X^{-1}$。若 $(X_0,\allowbreak F_0)$ 与 $(X_1,\allowbreak F_1)$ 为两个这样的对象，则态射群 $\Hom_{A^{\mathrm{strict}}}(X_0,\allowbreak X_1)=\Hom_A(X_0,\allowbreak X_1)$ 带自然对合 $E^{\mathrm{strict}}(\Phi)=F_1^{-1}\circ E(\Phi)\circ F_0$。\znote{201}{把协调作用换成严格作用的手法在范畴论中很常见：让对象额外携带“实现同构”的数据（这里是 $F_\pm$），态射则忘掉这些数据。于是新模式与原范畴等价，而群作用成为严格函子作用，不再需要额外的协调同构；代价是每个对象的严格化不唯一（相差一个自动同构），但对同调范畴无影响。第 (10b)、(12i) 小节用过同样的构造。[来源：原书 §(10b)、§(12i)。]}

转到几何情形：$M$ 为带角正合辛流形，带二次复体积形式 $\eta_M^2$。假设有保持全部给定结构的对合 $\iota_M$；这里包括 $\omega_M,\allowbreak I_M$ 及一形式 $\theta_M$，以后谈到正合辛流形的对合时，均以此为默认假设。添加
\SourcePage{210}
平凡分次，即零分次，以及线丛 $\beta=\xi_M$，将它变成膜自同构 $\iota_M^\#$。由一般理论，这定义 Donaldson–Fukaya 范畴的唯一自等价，作为我们的 $E$。此前考虑的是边界附近等于恒等的自同构；不过，除了正合性外，只需较弱的条件，即自同构在边界附近保持 $I_M$，而这里显然满足。用\stref{笔记}{11.18}的同构 $\xi_M|L\cong\lambda^{\mathrm{top}}(TL)$，可将 $E$ 对对象的作用写为
\begin{equation}\label{eq:14.7}
(L,\alpha^\#,P^\#)\longmapsto
(\iota_M(L),\alpha^\#\circ\iota_M,\iota_M^*(P^\#\otimes\lambda^{\mathrm{top}}(TL))).
\end{equation}
因为 $\iota_M^\#\circ\iota_M^\#=(\Id_M,\allowbreak 0,\allowbreak \xi_M\otimes\iota_M^*\xi_M)$ 平凡，$E^2=E\circ E$ 同构于恒等函子。更确切地，$\xi_M\otimes\iota_M^*\xi_M$ 的平凡化给出同构 $T:E^2\to\Id$。通过同构 $\lambda^{\mathrm{top}}(D\iota_M):\xi_M\to\iota_M^*\xi_M$ 有一个显然的平凡化，但我们要将它乘以 $-1$，并取相应的 $T$。唯一非平凡的协调条件\eref{14.6}归结为简单的事实：这一平凡化本身 $\Z/2$ 不变，其中作用是沿 $\iota_M$ 拉回并交换两个因子。

现在显然，预等变膜的定义正是上述抽象代数构造的一个实例。特别地，有 $\K$ 线性分次范畴 $H\F(M)^{\mathrm{strict}}$，对象为预等变膜，态射空间为 $HF^*(L_{0,\allowbreak +}^\#,\allowbreak L_{1,\allowbreak +}^\#)$。它等价于 $H\F(M)$，其唯一优势是带严格的 $\Z/2$ 作用。对对象，该作用就是之前描述的对合 $\epsilon$；对态射，规定
\begin{equation}\label{eq:14.8}
\begin{aligned}
\epsilon:HF^*(L_{0,+}^\#,L_{1,+}^\#)
&\xrightarrow{\cong}HF^*(L_{0,+}^\#\otimes\lambda^{\mathrm{top}}(TL_{0,+}),L_{1,+}^\#\otimes\lambda^{\mathrm{top}}(TL_{1,+}))\\
&\xrightarrow{\cong}HF^*(\iota_M(L_{0,+}^\#\otimes\lambda^{\mathrm{top}}(TL_{0,+})),
\iota_M(L_{1,+}^\#\otimes\lambda^{\mathrm{top}}(TL_{1,+})))\\
&\xrightarrow{\cong}HF^*(L_{0,-}^\#,L_{1,-}^\#).
\end{aligned}
\end{equation}
第一个同构使用 $\lambda^{\mathrm{top}}(TL_{k,\allowbreak +})$ 与同一个线丛 $\xi_M$ 的限制典范同构这一事实，比较\eref{14.7}；第二个来自分次辛自同构下的函子性；第三个来自与 $f_{0,\allowbreak -},\allowbreak f_{1,\allowbreak -}$ 的复合。

关于 $\Z/2$ 在 $H\F(M)^{\mathrm{strict}}$ 上的作用，还可陈述一些性质，基本上是\secref{12e}诸性质的等变版本。证明直接，故省略。
\begin{itemize}
\item \emph{改变符号。}令 $\beta\to M$ 为带非平凡 $\Z/2$ 作用的平凡线丛。显然，对所有预等变膜 $L_{k,\allowbreak +}^\#$，均有 $HF^*(L_{0,\allowbreak +}^\#,\allowbreak L_{1,\allowbreak +}^\#\otimes\beta|L_{1,\allowbreak +})=HF^*(L_{0,\allowbreak +}^\#,\allowbreak L_{1,\allowbreak +}^\#)$，但这些同构与\eref{14.8}反交换。特别地，若 $L_0^\#,\allowbreak L_1^\#$ 是真正的等变膜，则将其中一个与 $\beta$ 作张量积，会反转 $\Z/2$ 在 $HF^*(L_0^\#,\allowbreak L_1^\#)$ 上作用的符号，因而
\SourcePage{211}
在底域满足 $\operatorname{char}(\K)\ne2$ 时，交换不变部分与反不变部分。\znote{202}{把膜与 $\beta$ 张量积，相当于把所有 $f_\pm$ 同时乘以 $-e$。由于 $f_\pm$ 正是联系主部 $L_+$ 与影子 $L_-$ 的数据，这一操作在链层次给出同构，但它与 $\epsilon$ 的作用反交换，故在上同调上把 $\Z/2$ 作用变为负作用；特征不等于二时，正、负本征空间随之互换。这也说明 $\beta$ 的 $\Z/2$ 作用是本质的，不能通过重新选择平凡化消去。从表示论看，这是 $\Z/2$ 的一维实表示只有两个、而张量积给出群同态这一最简单事实的直接体现。[来源：原书 §(14a)，印刷页 201。]}
\item \emph{PSS 同构。}若 $L^\#$ 为等变膜，则 PSS 同构 $HF^*(L^\#,\allowbreak L^\#)\cong H^*(L;\K)$ 与两侧的 $\Z/2$ 作用相容。
\item \emph{同痕不变性。}取正合预等变 Lagrangian 膜的同痕，即族 $\{(L_{s,\allowbreak +}^\#,\allowbreak P_{s,\allowbreak -}^\#,\allowbreak f_{s,\allowbreak +},\allowbreak f_{s,\allowbreak -})\}$，$0\leq s\leq1$。如通常一样，$\{L_{s,\allowbreak +}^\#\}$ 确定 $HF^0(L_{0,\allowbreak +}^\#,\allowbreak L_{1,\allowbreak +}^\#)$ 的一个元素，其影子确定 $HF^0(L_{0,\allowbreak -}^\#,\allowbreak L_{1,\allowbreak -}^\#)$ 的一个元素。这两个元素实际上在 $\Z/2$ 作用 $\epsilon$ 下互为像。特别地，若两个正合等变膜在此类中同痕，则它们之间在 $H\F(M)^{\mathrm{strict}}$ 中存在不变同构。
\end{itemize}
